% Chapter 2 — The Picture: Map, Then Fold. Author: KLEPPMANN. Budget 5 pp.
\section{The Picture: Map, Then Fold}\label{ch:picture}

This chapter discharges constitution clauses C-3.1--C-3.11 (the map-then-fold picture, the three arrows, and the proposal-versus-authority split).

Everything this framework does is a \emph{map} followed by a \emph{fold}. The
constitution defines both words and draws the whole picture (C-3.1); this chapter does
not re-derive them --- it instantiates the picture on the six threads. In one line: the
world presents a list of events, $E = [e]$, and the ledger is the fold of that list ---
one event at a time, each step mapping the event to its transaction $tx$ against the
ledger built so far, then folding that $tx$ in, so the transactions $T = [tx]$ are output
the fold produces and never a list handed to it; and every number the ledger reports ---
a balance, a valuation, a profit and loss, a report --- is a further fold of that same
record. Four terms recur, each defined in Chapter~\ref{ch:objects} and used here
only at a gloss: a \textbf{unit} is anything that can be held or owed; a \textbf{wallet}
is where units are held; a \textbf{transaction} is the atomic bundle of changes the
ledger records; and a \textbf{projection} is any view computed from the record.

\subsection{The cast: six threads}\label{sec:cast}

Six instruments recur through this specification as threads. This section is their cast:
their terms, parties, and quantities are fixed here once and never vary, so a number met in
one chapter is the same number met in another. Every later episode opens with a pointer back
to this table and states only what that episode adds; a standing term is referred to here,
never re-established. This chapter follows one of them --- the one-touch OT-1 --- through the
whole pipeline; the rest are worked where they bite.

\medskip
{\small
\noindent\begin{tabular}{@{}lp{4.1cm}p{9.4cm}@{}}
\toprule
 & Instrument, kind, parties & Fixed terms and headline numbers \\
\midrule
T1 & FUT-IDX --- cash-settled future on IDX; W-ALPHA / the CCP & multiplier 10, USD,
  settled-to-market; buy 1 at 995.00; settlement prints 1{,}000.00 / 990.00 /
  1{,}005.00; cumulative VM \USD{100} $=$ PnL \\
T2 & OT-1 --- one-touch on stock OMEGA; W-BANK writes, W-ALPHA holds & barrier
  breached when OMEGA's official close is at or below 80.00 --- a discretely-monitored
  (closing-price) one-touch; payout \USD{1{,}000{,}000}; spot 95.00 at inception; marked
  \USD{400{,}000} while live; knock at close 79.00 \\
T3 & ACME --- cash dividend; issuer / W-ALPHA & 1.50 per share; W-ALPHA holds 10{,}000
  shares; entitlement \USD{15{,}000} \\
T4 & ACME --- 2-for-1 split; issuer / W-ALPHA & 10{,}000 $\to$ 20{,}000 shares; price
  frame 100.00 $\to$ 50.00; declared dividend figure 1.50 $\to$ 0.75 \\
T5 & VS-1 --- OTC variance swap on IDX; bilateral & one year, 252 fixings, strike 400
  variance points, \USD{1{,}000} per point; realised variance 441; settlement
  \USD{41{,}000} \\
T6 & TRS-1 --- total return swap; reference W-QIS in virtual ledger V-1 (strategy S-QIS) &
  notional \USD{10{,}000{,}000}; financing 4\% p.a., quarterly; first reset NAV
  10{,}300{,}000; net move \USD{200{,}000} \\
\bottomrule
\end{tabular}
}
\medskip

\subsection{The world presents events}

The world presents the system with a list of events. Some arrive from outside: trades
executed in the market, corporate-action notices, fixings and other observations,
captured and recorded at the boundary. Others are not received but noticed. Every unit
declares, from its terms, the conditions its life requires watching --- its
\textbf{watches} (Chapter~\ref{ch:objects}) --- and the Event Monitor evaluates those
watches against the clock and the recorded observations, emitting an event onto the
record the moment a condition is met. The future FUT-IDX declares such watches at
registration: one for each daily settlement print, one for its expiry; the one-touch OT-1
declares one, its underlying's recorded close against the barrier. A coupon date
arriving, an expiry falling due, a close printing below a barrier --- each becomes an
event the same way a trade does, and once emitted the two kinds are indistinguishable:
both are on the record.

\subsection{The map: each event becomes a transaction}

Each event is mapped to a transaction, and the responsibility for the map divides in two.
The \textbf{smart contracts} are the mapping functions: each is a stateless piece of code
that, given one event and the current state of the ledger, computes the transaction the
event requires (Chapter~\ref{ch:contracts}). The \textbf{Events Executor} fires the map
one event at a time as the fold advances: for each event, in the total order, it fires the
responsible contract, hands it the event together with the current projection of the
ledger, and collects the proposed transaction $tx$ --- so $T=[tx]$ is built by the fold,
not consumed by it. Many contracts, one Events Executor; the contracts carry all the
product knowledge, the Events Executor carries none (Chapter~\ref{ch:machines}).

The map is not blind to state. To turn a dividend event into cash moves, the contract
must read who holds the stock at that moment --- the fold of everything admitted so far.
Map and fold therefore advance together, one event at a time, in a single total order:
the transaction for event $e$ is computed against the ledger that the transactions for
all earlier events have already built.

\subsection{The fold: the transactions become the ledger}

There is no list of transactions waiting to be folded --- the fold builds it. As each
proposed transaction arrives in the total order, the \textbf{Transaction Executor} applies
it to the state of the system: balances, and the three homes of unit and position state
(Chapter~6, \emph{State: The Three Homes}). It serialises, validates, and applies one
transaction at a time, $(\,tx, \mathrm{Ledger}) \to \mathrm{Ledger}$, and the immutable log
is precisely the list the fold has so far built --- the fold's output, never its input.
Nothing else writes. And because the log is one list, every other view is a further fold of it: a
balance sums the moves that touch a wallet, a valuation folds the composition against a
price vector (Chapter~\ref{ch:valuation}), a report folds it into a presentation
(Chapter~\ref{ch:reporting}). One record, and many folds of it --- never many records.

\subsection{Two axes of time: what was known, and what is in force}\label{sec:bitemporal}

A correction arrives late. IDX's official close for day~2 is $990.00$, but the print
reaches the boundary only on day~4 --- after day~3 has already settled against the
prints then on the record. The system now holds two honest questions with different
answers. \emph{As the book stood at the close of day~3}, day~2's variation margin was
computed without the $990.00$: that is what was reported, and no later arrival may rewrite
it. \emph{As the book stands once the correction is in}, day~2's margin is recomputed with
the $990.00$ in force and the tail from day~2 forward re-derived. Both answers are correct.
They answer different questions.

Two coordinates, not one, therefore fix a historical view. The first is the \emph{log
prefix}: how far down the immutable log the reconstruction reads --- what had been recorded
by a chosen point. This is \emph{transaction time}, the axis of what was known. The second
is the \emph{knowledge horizon}: which of the recorded observations are treated as in
force. This is \emph{valid time}, the axis of what a recorded observation says about the
world, apart from when it was recorded. A back-dated observation --- day~2's close, arriving
on day~4 --- sits at valid-time index $k=2$ on the second axis while it sits late on the
first. The projection reads both:

\begin{lstlisting}
project :: LogPrefix -> KnowledgeHorizon -> View
-- LogPrefix        : how far down the immutable log to read   (transaction time)
-- KnowledgeHorizon : which recorded observations are in force (valid time)
\end{lstlisting}

Fix both coordinates and the view is a deterministic fold, computed the same by any party.
Two facts follow, and they are the whole of bitemporality. First, the \emph{as-known-at}
view --- the projection over a fixed prefix --- is never rewritten by a later correction: a
correction is a new recorded observation appended further down the log, so it lies outside
the fixed prefix and cannot reach into it. Second, the \emph{as-of} view --- the projection
over the prefix that already includes the correction --- recomputes bit-exactly from that
prefix, because the fold reads only recorded observations and is deterministic. The late
close does not edit day~2: it is recorded at its own valid-time index~$k=2$, the fold from
that index forward is re-derived, and the earlier as-known-at view survives beside it.
Theorem~\ref{thm:timetravel} states this over the pair of axes and proves both facts.

The constitution names one axis. \S 12 reconstructs a historical state ``by replaying the
immutable transaction log up to the desired timestamp'' --- one timestamp, one axis --- yet
a back-dated observation forces the two just drawn. This is a genuine conflict between a
single-axis guarantee and a system that has needed both axes since the market-data boundary
of Chapter~\ref{ch:marketdata}. It is parked, with the exact proposed amendment text, as
entry PARK-2 of the parking index (Chapter~\ref{ch:scope}); the machinery this section
states is what that amendment names.

\subsection{Data, functions, and machines}

Three kinds of thing appear in this picture, and they are never confused: data,
functions, and machines. The data kinds are Event, Transaction, and Ledger. The functions are
the arrows between the data. The machines are what evaluate the arrows --- one machine per
arrow, and they are the subject of Chapter~\ref{ch:machines}. Typed:

\begin{lstlisting}
watch    : (Record, Clock)       -> [Event]      -- the Event Monitor
contract : (Event, Ledger)       -> Transaction  -- the Events Executor
apply    : (Transaction, Ledger) -> Ledger       -- the Transaction Executor

step (ledger, e) = apply (contract (e, ledger), ledger)  -- map, then fold
final ledger     = fold step over E = [e]
\end{lstlisting}

The three arrows differ in what they read. \texttt{watch} reads the Record and the clock
and yields events; \texttt{contract} and \texttt{apply} consume the Ledger --- the folded
state --- and never the clock. One \texttt{step} is a map (fire the contract) immediately
followed by a fold (apply the result), and the whole run is that step folded over
$E=[e]$. Map, then fold: the phrase is the architecture, not a metaphor for it.

An observation is not outside this picture. It enters as a moveless
\texttt{Transaction} through \texttt{apply} --- a transaction with an empty move list
that folds the observation into a home (Chapter~\ref{ch:marketdata}) --- so the Record
the \texttt{watch} arrow reads and the Ledger \texttt{contract} and \texttt{apply}
consume are two views of one log, never two stores. The observations a \texttt{contract}
reads are therefore home facts already folded into the Ledger, and the single writer of
\texttt{apply} is the only writer of any observation.

\subsection{The clock, and proposal versus authority}

Two properties of the picture carry most of its weight.

\emph{The clock is the single input that is not on the record}, and the Event Monitor is
the single place it is consulted. The event the Monitor emits is itself recorded, so
everything downstream of the emission is a function of the record alone --- which is
exactly why a late firing produces the identical transaction, and why, given the record
and the clock, any party can recompute which events should have been emitted, and when
(Chapter~\ref{ch:machines}).

\emph{The map is proposed; the fold is authoritative.} A contract's output is a proposal,
not a fact: a transaction joins the list only when the Transaction Executor admits it. The
smart contracts, the Event Monitor, and the Events Executor sit outside the trust boundary
and may only propose; the one Transaction Executor writes. Either arrow may refuse --- a
contract may find nothing due; the door rejects a malformed transaction and writes nothing
--- and a refusal is a returned value, never a half-applied state. Because both the events
and the transactions are on the record, economic correctness is checkable after the fact:
re-run the map on the recorded event and compare the result against the recorded
transaction (Chapter~15, \emph{Testability and the Executable-Check Regime}). Verification
lives with the events; authority lives with the transactions.

\subsection{One instrument walks the pipeline}

OT-1 (Cast, \S\ref{sec:cast}) walks the pipeline whole. Its standing terms are fixed above;
this section adds only the walk across the three arrows. At registration its terms declared
one watch: OMEGA's recorded official closes against the barrier.

\textbf{Trigger (watch).} OMEGA's official close of 79.00 is recorded. The declared
condition is met, and the Event Monitor emits the touch event
onto the record. No ledger state has changed: the Monitor emits events, and an event is an
input to the map, not a write.

\textbf{Contract (map).} The Events Executor fires OT-1's contract, handing it the touch
event and the current projection of the ledger. The contract reads the terms, the holder
(\mbox{W-ALPHA}), and the unit's state (\emph{live}), and
returns one transaction:

\begin{center}
\begin{tabular}{@{}rll@{}}
\toprule
 & Kind & Content \\
\midrule
1 & unit-state change & OT-1: UnitStatus \emph{live} $\to$ \emph{triggered} \\
2 & move & owned(USD) 1{,}000{,}000: W-BANK $\to$ W-ALPHA \\
\bottomrule
\end{tabular}
\end{center}

\textbf{Apply (fold).} The Transaction Executor admits it --- the move's paired legs leave
conservation undisturbed, the cause-derived identifier names the touch event, and the
writer is OT-1's own contract --- and the fold advances: the ledger now shows the cash with
\mbox{W-ALPHA} and the option \emph{triggered}. Where OT-1 stands encumbered under a
collateral agreement the payment leg routes differently, and Chapter~\ref{ch:collateral}
works that same knock start to finish; the pipeline is the one drawn here either way.

\emph{One list of events, one map, one fold: the option's whole knock is a single event
carried across the three arrows, and nothing off the record ever touches it.}

\subsection{Not event sourcing}

The resemblance to event sourcing is real, and the difference matters. The Ledger is
event-driven because finance is: trades execute, prices print, coupons mature, and
barriers are touched whether or not the software cares, so events are recorded because
reality presents them, not because a pattern was chosen. The difference is where authority
sits. In event sourcing the event stream is the canonical history, state is whatever the
current code derives from it, and reinterpreting history through new code is a feature.
Here the admitted transaction log alone is the canonical history, and two properties set it
apart. First, there is a single total order, not a set of per-aggregate streams stitched
back together: every view is a projection of that one log, and two projections of it cannot
disagree. Second, replay re-reads transactions and never re-executes contracts
(Chapter~\ref{ch:machines}), so history cannot be silently reinterpreted. Events are inputs
to the map, never the accounting record: contracts transform them into proposed
transactions, and the Transaction Executor admits or refuses. Correcting the record is
neither reinterpretation nor an edit, and the two axes of Section~\ref{sec:bitemporal} say
why: a wrong interpretation is repaired by a compensating transaction in the open
(Chapter~15), and a corrected observation --- a close that arrives late or is later restated
--- is a new recorded observation at its own valid-time index, appended on the
transaction-time axis rather than written over the old one. The as-known-at view over any
earlier prefix is therefore preserved, and the as-of view recomputes from the prefix that
carries the correction; neither overwrites the other. The events stay on the record for one
purpose --- so the map can be re-run and verified.

\begin{secondtelling}
A fold with a fixed step is a \emph{catamorphism}: the unique map out of the list $E=[e]$
determined by one rule for the empty list and one rule for prepending an element --- here
\texttt{step}. Because the step is fixed and the list carries its own order, the ledger is
that unique image, and every projection --- balance, valuation, report --- is another
catamorphism over the same list. ``Map, then fold'' names the composite; that the fold is
a catamorphism is why one record admits many views and no two of them can disagree.
\end{secondtelling}
